\subsection{The finite components over the rational function field}

Put $K=\mathbb Q(T)$ and let $\Omega_h$ be the closed points of
$\mathbb A^1_{\mathbb Q}$. We equip $\Omega_h$ with its discrete
$\sigma$-algebra and counting measure, as in
\cite[\S3.2.5]{chen-moriwaki}.

\begin{definition}[The horizontal absolute values]
\label{def:rational-horizontal-place}
\source{chen-moriwaki-adelic:page-0188}
\lean{ArakelovGeometry.AdelicCurve.RationalFunction.HorizontalPlace,ArakelovGeometry.AdelicCurve.RationalFunction.horizontalMeasurableSpace,ArakelovGeometry.AdelicCurve.RationalFunction.horizontalPolynomial,ArakelovGeometry.AdelicCurve.RationalFunction.horizontalPolynomial_isPrimitive,ArakelovGeometry.AdelicCurve.RationalFunction.horizontalPolynomial_ne_zero,ArakelovGeometry.AdelicCurve.RationalFunction.horizontalPolynomial_span,ArakelovGeometry.AdelicCurve.RationalFunction.horizontalHeight,ArakelovGeometry.AdelicCurve.RationalFunction.one_le_horizontalHeight,ArakelovGeometry.AdelicCurve.RationalFunction.horizontalAbv,ArakelovGeometry.AdelicCurve.RationalFunction.horizontalAbv_apply,ArakelovGeometry.AdelicCurve.RationalFunction.horizontalAbv_rat}
\leanok
For $x\in\Omega_h$, choose a generator of the corresponding nonzero prime
ideal of $\mathbb Q[T]$, clear its denominators, and take its primitive
part $F_x\in\mathbb Z[T]$. Then $F_x\ne0$ and generates the same ideal
over $\mathbb Q$. Equip $\Omega_h$ with the discrete $\sigma$-algebra.
Let $H(x)=M(F_x)$ be its Mahler measure. We have $H(x)\geq1$.
The horizontal absolute value is
\[
 |0|_x=0,\qquad |f|_x=H(x)^{-\operatorname{ord}_x(f)}\quad(f\ne0).
\]
It restricts to the trivial absolute value on $\mathbb Q$; when
$H(x)=1$, it is the trivial absolute value on $K$ itself.
\end{definition}
\begin{proof}
Clearing denominators multiplies a generator by a nonzero integer.
Taking the primitive part divides it by its nonzero integral content.
Both scalars are units in $\mathbb Q[T]$, so the resulting polynomial
still generates the ideal of $x$, is nonzero, and is primitive.
For every nonzero integral polynomial the Mahler measure is at least
the absolute value of its leading coefficient, hence at least one.
If $H(x)>1$, the valuation inequalities and multiplicativity of
$\operatorname{ord}_x$ imply the nonarchimedean triangle inequality and
multiplicativity of the displayed absolute value. If $H(x)=1$, the
displayed formula is the trivial absolute value and satisfies the same
axioms directly. Finally a nonzero constant is a unit in
$\mathbb Q[T]$, so its order at $x$ is zero and its absolute value is one.
\end{proof}

\begin{definition}[The prime Gauss absolute values]
\label{def:rational-prime-place}
\source{chen-moriwaki-adelic:page-0188,chen-moriwaki-adelic:page-0189}
\lean{ArakelovGeometry.AdelicCurve.RationalFunction.rationalPadicAbv,ArakelovGeometry.AdelicCurve.RationalFunction.primeMeasurableSpace,ArakelovGeometry.AdelicCurve.RationalFunction.isNonarchimedean_rationalPadicAbv,ArakelovGeometry.AdelicCurve.RationalFunction.primePolynomialValuation,ArakelovGeometry.AdelicCurve.RationalFunction.primePolynomialValuation_ne_zero,ArakelovGeometry.AdelicCurve.RationalFunction.primeValuation,ArakelovGeometry.AdelicCurve.RationalFunction.primeAbv,ArakelovGeometry.AdelicCurve.RationalFunction.primeAbv_polynomial,ArakelovGeometry.AdelicCurve.RationalFunction.isNonarchimedean_primeAbv,ArakelovGeometry.AdelicCurve.RationalFunction.primeAbv_num_div_denom,ArakelovGeometry.AdelicCurve.RationalFunction.primeAbv_rat}
\leanok
For a rational prime $p$ and a polynomial $P=\sum_i a_iT^i$, put
\[
 \lVert P\rVert_p=\max_i|a_i|_p,
\]
where the maximum is zero when $P=0$ and $|\cdot|_p$ is the usual
nonarchimedean rational $p$-adic absolute value. Equip the prime index
set with the discrete $\sigma$-algebra. The Gauss norm, regarded as a
valuation with values in $\mathbb R_{\geq0}$, is nonzero on every
nonzero polynomial. It defines a nonarchimedean absolute value on
$\mathbb Q[T]$ and extends to $K$ by
\[
 |P/Q|_p=\lVert P\rVert_p/\lVert Q\rVert_p\qquad(Q\ne0).
\]
Its restriction to the constant field is the original rational
$p$-adic absolute value.
\end{definition}
\begin{proof}
The coefficientwise nonarchimedean inequality gives
$\lVert P+Q\rVert_p\leq\max(\lVert P\rVert_p,\lVert Q\rVert_p)$;
nonnegativity and vanishing only on the zero polynomial are immediate.
For nonzero $P,Q$, choose the least indices $i,j$ attaining their
respective coefficient maxima. In the coefficient of degree $i+j$ of
$PQ$, the term $a_ib_j$ has absolute value
$\lVert P\rVert_p\lVert Q\rVert_p$. Every other term has strictly
smaller absolute value: a smaller first index gives a smaller first
factor, and a larger first index gives a second index smaller than $j$.
The strong triangle inequality then shows that this coefficient has
absolute value exactly $\lVert P\rVert_p\lVert Q\rVert_p$.
All coefficients are bounded above by that product, so the Gauss norm
is multiplicative. The zero cases follow immediately.
If $P/Q=P'/Q'$ with nonzero denominators, then $PQ'=P'Q$;
multiplicativity and division by the nonzero denominator norms prove
that the quotient formula is independent of the presentation.
Applying this formula after putting two fractions over a common
denominator proves the nonarchimedean inequality on $K$.
It also proves multiplicativity and vanishing precisely at zero.
For a constant polynomial the Gauss norm is its single coefficient's
absolute value, which proves the restriction assertion.
\end{proof}

\begin{lemma}[Integrability at the rational primes]
\label{lem:rational-prime-log-integrable}
\source{chen-moriwaki-adelic:page-0188,chen-moriwaki-adelic:page-0189}
\uses{def:rational-prime-place}
\lean{ArakelovGeometry.AdelicCurve.RationalFunction.rationalPadicAbv_eq_one_of_large,ArakelovGeometry.AdelicCurve.RationalFunction.finite_rationalPadicAbv_ne_one,ArakelovGeometry.AdelicCurve.RationalFunction.finite_primeAbv_polynomial_ne_one,ArakelovGeometry.AdelicCurve.RationalFunction.finite_support_log_primeAbv,ArakelovGeometry.AdelicCurve.RationalFunction.integrable_log_primeAbv,ArakelovGeometry.AdelicCurve.RationalFunction.primeAdelicStructure}
\leanok
For $a=m/n\in\mathbb Q^\times$ in lowest terms, $|a|_p=1$ whenever
$p>\max(|m|,n)$, so only finitely many primes have $|a|_p\ne1$.
For every nonzero polynomial $P\in\mathbb Q[T]$, only finitely many
primes have $\lVert P\rVert_p\ne1$.
For $f\in K^\times$, the function $p\mapsto\log|f|_p$ has finite
support on the set of rational primes. It is therefore integrable for
counting measure and defines the prime adelic structure on $K$.
\end{lemma}
\begin{proof}
For $a=m/n\in\mathbb Q^\times$ in lowest terms, a prime exceeding
$\max(|m|,n)$ divides neither $m$ nor $n$, so $|a|_p=1$.
There are only finitely many smaller primes. Thus each nonzero rational
coefficient has absolute value one outside a finite set of primes.
For a nonzero polynomial, take the union of these exceptional sets
over its finitely many nonzero coefficients. Outside this finite
union all nonzero coefficients have norm one, so their maximum is one.
Finally write $f=P/Q$ with $P,Q\ne0$. Outside the union of the finite
exceptional sets for $P$ and $Q$, its absolute value is $1/1=1$.
The logarithm vanishes there. On the discrete prime space measurability
is automatic, and counting measure integrates its absolute value as a
finite sum. This proves all the adelic integrability axioms.
\end{proof}

\begin{lemma}[Integrability at the horizontal places]
\label{lem:rational-horizontal-log-integrable}
\source{chen-moriwaki-adelic:page-0188,chen-moriwaki-adelic:page-0189}
\uses{def:rational-horizontal-place}
\lean{ArakelovGeometry.AdelicCurve.RationalFunction.finite_support_log_horizontalAbv,ArakelovGeometry.AdelicCurve.RationalFunction.integrable_log_horizontalAbv,ArakelovGeometry.AdelicCurve.RationalFunction.horizontalAdelicStructure}
\leanok
For $f\in K^\times$, the function $x\mapsto\log|f|_x$ has finite
support on $\Omega_h$. Consequently it is integrable for counting
measure and defines the horizontal adelic structure on $K$.
\end{lemma}
\begin{proof}
Write $f=P/Q$ with nonzero $P,Q\in\mathbb Q[T]$.
Outside the finite union of the prime divisors of $P$ and $Q$,
$\operatorname{ord}_x(f)=0$. The preceding formula then gives
$|f|_x=1$, so its logarithm is zero.
Every function on the discrete space $\Omega_h$ is measurable.
Since each singleton has counting measure one, the integral of the
absolute value of a finitely supported real function is a finite sum
of finite real numbers. This proves integrability, and hence every
axiom of the asserted adelic structure.
\end{proof}
